\section{AIME I 1992}
        \begin{enumerate}
        	\item (\textit{AIME I 1992}) Find the sum of all positive rational numbers that are less than 10 and that have denominator 30 when written in lowest terms.


 

	\item (\textit{AIME I 1992}) A positive integer is called ascending if, in its decimal representation, there are at least two digits and each digit is less than any digit to its right. 
How many ascending positive integers are there?


 

	\item (\textit{AIME I 1992}) A tennis player computes her win ratio by dividing the number of matches she has won by the total number of matches she has played. At the start of a weekend, her win ratio is exactly $0.500$. During the weekend, she plays four matches, winning three and losing one. At the end of the weekend, her win ratio is greater than $.503$. What's the largest number of matches she could've won before the weekend began?


 

	\item (\textit{AIME I 1992}) In Pascal's Triangle, each entry is the sum of the two entries above it. The first few rows of the triangle are shown below.


 
 \[\begin{array}{c@{\hspace{8em}} c@{\hspace{6pt}}c@{\hspace{6pt}}c@{\hspace{6pt}}c@{\hspace{4pt}}c@{\hspace{2pt}} c@{\hspace{2pt}}c@{\hspace{2pt}}c@{\hspace{2pt}}c@{\hspace{3pt}}c@{\hspace{6pt}} c@{\hspace{6pt}}c@{\hspace{6pt}}c} \vspace{4pt} \text{Row 0: } &    &    &     &     &    &    & 1 &     &     &    &    &    &  \\\vspace{4pt} \text{Row 1: } &    &    &     &     &    & 1 &    & 1  &     &    &    &    &  \\\vspace{4pt} \text{Row 2: } &    &    &     &     & 1 &    & 2 &     & 1  &    &    &    &  \\\vspace{4pt} \text{Row 3: } &    &    &     &  1 &    & 3 &    & 3  &     & 1 &    &    &  \\\vspace{4pt} \text{Row 4: } &    &    & 1  &     & 4 &    & 6 &     & 4  &    & 1 &    &  \\\vspace{4pt} \text{Row 5: } &    & 1 &     & 5  &    &10&    &10 &     & 5 &    & 1 &  \\\vspace{4pt} \text{Row 6: } & 1 &    & 6  &     &15&    &20&     &15 &    & 6 &    & 1 \end{array}\]In which row of Pascal's Triangle do three consecutive entries occur that are in the ratio $3: 4: 5$? 


 

	\item (\textit{AIME I 1992}) Let $S^{}_{}$ be the set of all rational numbers $r^{}_{}$, $0^{}_{}<r<1$, that have a repeating decimal expansion in the form $0.abcabcabc\ldots=0.\overline{abc}$, where the digits $a^{}_{}$, $b^{}_{}$, and $c^{}_{}$ are not necessarily distinct. To write the elements of $S^{}_{}$ as fractions in lowest terms, how many different numerators are required?


 

	\item (\textit{AIME I 1992}) For how many pairs of consecutive integers in $\{1000,1001,1002^{}_{},\ldots,2000\}$ is no carrying required when the two integers are added?


 

	\item (\textit{AIME I 1992}) Faces $ABC^{}_{}$ and $BCD^{}_{}$ of tetrahedron $ABCD^{}_{}$ meet at an angle of $30^\circ$. The area of face $ABC^{}_{}$ is $120^{}_{}$, the area of face $BCD^{}_{}$ is $80^{}_{}$, and $BC=10^{}_{}$. Find the volume of the tetrahedron.


 

	\item (\textit{AIME I 1992}) For any sequence of real numbers $A=(a_1,a_2,a_3,\ldots)$, define $\Delta A^{}_{}$ to be the sequence $(a_2-a_1,a_3-a_2,a_4-a_3,\ldots)$, whose $n^{th}$ term is $a_{n+1}-a_n^{}$. Suppose that all of the terms of the sequence $\Delta(\Delta A^{}_{})$ are $1^{}_{}$, and that $a_{19}=a_{92}^{}=0$. Find $a_1^{}$.


 

	\item (\textit{AIME I 1992}) Trapezoid $ABCD^{}_{}$ has sides $AB=92^{}_{}$, $BC=50^{}_{}$, $CD=19^{}_{}$, and $AD=70^{}_{}$, with $AB^{}_{}$ parallel to $CD^{}_{}$. A circle with center $P^{}_{}$ on $AB^{}_{}$ is drawn tangent to $BC^{}_{}$ and $AD^{}_{}$. Given that $AP^{}_{}=\frac mn$, where $m^{}_{}$ and $n^{}_{}$ are relatively prime positive integers, find $m+n^{}_{}$.


 

	\item (\textit{AIME I 1992}) Consider the region $A^{}_{}$ in the complex plane that consists of all points $z^{}_{}$ such that both $\frac{z^{}_{}}{40}$ and $\frac{40^{}_{}}{\overline{z}}$ have real and imaginary parts between $0^{}_{}$ and $1^{}_{}$, inclusive. What is the integer that is nearest the area of $A^{}_{}$?


 

	\item (\textit{AIME I 1992}) Lines $l_1^{}$ and $l_2^{}$ both pass through the origin and make first-quadrant angles of $\frac{\pi}{70}$ and $\frac{\pi}{54}$ radians, respectively, with the positive x-axis. For any line $l^{}_{}$, the transformation $R(l)^{}_{}$ produces another line as follows: $l^{}_{}$ is reflected in $l_1^{}$, and the resulting line is reflected in $l_2^{}$. Let $R^{(1)}(l)=R(l)^{}_{}$ and $R^{(n)}(l)^{}_{}=R\left(R^{(n-1)}(l)\right)$. Given that $l^{}_{}$ is the line $y=\frac{19}{92}x^{}_{}$, find the smallest positive integer $m^{}_{}$ for which $R^{(m)}(l)=l^{}_{}$.


 

	\item (\textit{AIME I 1992}) In a game of \textit{Chomp}, two players alternately take bites from a 5-by-7 grid of unit squares. To take a bite, a player chooses one of the remaining squares, then removes ("eats") all squares in the quadrant defined by the left edge (extended upward) and the lower edge (extended rightward) of the chosen square. For example, the bite determined by the shaded square in the diagram would remove the shaded square and the four squares marked by $\times.$ (The squares with two or more dotted edges have been removed from the original board in previous moves.) 


 
                    \begin{figure}[H]
                        \centering
                        \includegraphics[width=0.5\textwidth]{Problem/AIME_I_1992_12_Problem_Figure_1.png}
                    \end{figure}
                    


 The object of the game is to make one's opponent take the last bite. The diagram shows one of the many subsets of the set of 35 unit squares that can occur during the game of Chomp. How many different subsets are there in all? Include the full board and empty board in your count.


 

	\item (\textit{AIME I 1992}) Triangle $ABC^{}_{}$ has $AB=9^{}_{}$ and $BC: AC=40: 41^{}_{}$. What's the largest area that this triangle can have?


 

	\item (\textit{AIME I 1992}) In triangle $ABC^{}_{}$, $A'$, $B'$, and $C'$ are on the sides $BC$, $AC^{}_{}$, and $AB^{}_{}$, respectively. Given that $AA'$, $BB'$, and $CC'$ are concurrent at the point $O^{}_{}$, and that $\frac{AO^{}_{}}{OA'}+\frac{BO}{OB'}+\frac{CO}{OC'}=92$, find $\frac{AO}{OA'}\cdot \frac{BO}{OB'}\cdot \frac{CO}{OC'}$.


 

	\item (\textit{AIME I 1992}) Define a positive integer $n^{}_{}$ to be a factorial tail if there is some positive integer $m^{}_{}$ such that the decimal representation of $m!$ ends with exactly $n$ zeroes. How many positive integers less than $1992$ are not factorial tails?


 

\end{enumerate}